\exercisesection{1}

\begin{enumerate}
\def\labelenumi{\arabic{enumi}.}
\tightlist
\item
  Let A be a graded commutative ring of type Z and \((A_{n})\) its graduation; we set \(A^{\geq} = \bigoplus_{n \geq 0} A_{n}, A^{\leq} = \bigoplus_{n \leqslant 0} A_{n}\) ; these are graded subrings of A.

\begin{enumerate}
\def\labelenumii{(\alph{enumii})}
\item
  For every homogeneous element \(f\) of \(\mathbf{A}\) of degree \(d\) , the ring of fractions \(\mathbf{A}_f\) (corresponding to the multiplicative subset consisting of the \(f^n\) , where \(n \geqslant 0\) ) has a canonical graded ring structure (Algebra, Chapter II, §11); we denote by \(\mathbf{A}_{(f)}\) the subring of \(\mathbf{A}_f\) consisting of elements of degree 0. Show that, if \(d > 0\) , then \((\mathbf{A}^{\geq})_f = \mathbf{A}_f, \mathbf{A}_{(f)}\) is isomorphic to \(\mathbf{A}^{(d)} / (f - 1)\mathbf{A}^{(d)}\) and \(((\mathbf{A}_f)^{\geq})_{f/1}\) is a graded ring isomorphic to \(\mathbf{A}_f\) .
\item
  Let the polynomial ring \(\mathbf{B} = \mathbf{A}[\mathbf{X}] = \mathbf{A} \otimes_{\mathbf{Z}} \mathbf{Z}[\mathbf{X}]\) be given the graduation the tensor product of those on \(\mathbf{A}\) and \(\mathbf{Z}[\mathbf{X}]\) , a graduation compatible with the ring structure on \(\mathbf{B}\) . Show that, if \(d > 0\) , \(\mathrm{B}_{(f)}\) is isomorphic to \((\mathrm{A}_f)^{\leqslant}\) and \((\mathrm{A}^{(d)})_f\) is a graded ring isomorphic to \(\mathrm{A}_{(f)} \otimes_{\mathbf{Z}} \mathbf{Z}[\mathbf{X}, \mathbf{X}^{-1}]\) .
\item
  Let \(g\) be another homogeneous element of \(A\) of degree \(e\) . Show that, if \(d > 0\) and \(e > 0\) , \(A_{(fg)}\) is isomorphic to \((A_{(f)})_{g^d / f^e}\) .\\
\item[* (d)] Suppose \(d > 0\) and \(A_n = \{0\}\) for \(n < 0\) . Show that, if \(A\) is Noetherian, so is \(A_{(f)}\) (use §2, no. 10, Corollary 4 to Theorem 2). *
\end{enumerate}

\item
  Let A be a graded commutative ring of type Z such that \(A_{n}=0\) for n\textless0. For all \(n\geqslant0\) , let \(A_{[n]}\) denote the graded ideal \(\bigoplus_{m\geqslant n}A_{m}\) of A; the A-algebra \(A^{\sharp}=\bigoplus_{n\geqslant0}A_{[n]}\) is a ring graded by the \(A_{[n]}=A_{n}^{\sharp}\) ; for all \(f\in A_{d}(d>0)\) , let \(f^{\sharp}\) denote the element of \(A^{\sharp}\) whose component are zero except for that of degree d, which is equal to f.

\begin{enumerate}
\def\labelenumii{(\alph{enumii})}
\item
  Show that, if \(f \in \mathbf{A}_d\) , where \(d > 0\) , the ring \(\mathrm{A}_{(f^{\sharp})}^{\sharp}\) is isomorphic to \((\mathbf{A}_f)^{\geqslant}\) (notation of Exercise 1).
\item
  For \(\mathbf{A}^{\natural}\) to be a finitely generated A-algebra, it is necessary and sufficient that \(\mathbf{A}\) be a finitely generated \(\mathbf{A}_0\) -algebra.
\item
  In order that \(\mathbf{A}_{n+1}^{\natural} = \mathbf{A}_1^{\natural}\mathbf{A}_n^{\natural}\) (resp. \(\mathbf{A}_n^{\natural} = (\mathbf{A}_1^{\natural})^n\) ) for \(n \geqslant n_0\) , it is necessary and sufficient that \(\mathbf{A}_{n+1} = \mathbf{A}_1\mathbf{A}_n\) (resp. \(\mathbf{A}_n = (\mathbf{A}_1)^n\) ) for \(n \geqslant n_0\) .
\end{enumerate}

\item
  Let K be a field, B the polynomial ring \(\mathrm{K}[X,Y]\) with the graduation defined by the total degree, C the graded subring of B generated by X and \(Y^{2}\) and a the graded ideal of C generated by \(Y^{4}\) . Show that, in the graded ring A = C/a, \(A_{n+1} = A_{1}A_{n}\) for \(n \geqslant 2\) but \(A_{n} \neq (A_{1})^{n}\) for \(n \geqslant 6\) .
\end{enumerate}
% semantic-fragments: legacy-input-seam
\relax{}
